\documentclass[12pt]{article}

\usepackage[a4paper, margin=2.5cm]{geometry}
\usepackage{amsmath}
\usepackage{amssymb}
\usepackage{amsthm}
\usepackage[colorlinks=true, linkcolor=blue, citecolor=blue, urlcolor=blue]{hyperref}
\usepackage{booktabs}
\usepackage{natbib}

\newtheorem{theorem}{Theorem}
\newtheorem{proposition}[theorem]{Proposition}
\newtheorem{definition}[theorem]{Definition}
\newtheorem{remark}[theorem]{Remark}

\newcommand{\Mcl}{\mathcal{M}_{\mathrm{cl}}}
\newcommand{\Ccal}{\mathcal{C}}
\newcommand{\Fcal}{\mathcal{F}}
\newcommand{\Sadm}{\mathcal{S}_{\mathrm{adm}}}
\newcommand{\Scl}{\mathcal{S}_{\mathrm{cl}}}
\newcommand{\phig}{\varphi}

\title{Gravitational Analogy from Constraint-Manifold Curvature\\in $\phig$-Structured Geometry}
\author{%
  Lee Smart\\[4pt]
  \textit{Institute of Vibrational Field Dynamics}\\[2pt]
  \texttt{contact@vibrationalfielddynamics.org}\\[2pt]
  \texttt{@vfd\_org}%
}
\date{April 2026}

\begin{document}

\maketitle

\noindent\textbf{Status:} Preprint (not peer-reviewed)

\begin{abstract}
Paper~IX~\citep{SmartIX} constructed a continuous-limit formalism in which the closed state set is represented as a constraint manifold $\Mcl$ equipped with an induced Riemannian metric. The present paper develops the first gravity-facing interpretation of this framework: curvature of $\Mcl$ is proposed as a structural analogue of gravitational influence, and constrained motion on $\Mcl$ is interpreted as geodesic-like evolution.

Within this interpretation, free closure-compatible motion follows geodesics of the induced metric on $\Mcl$; curvature arising from constraint interaction deflects nearby trajectories in a manner structurally analogous to gravitational lensing; and regions of approximate local flatness provide a constraint-manifold analogue of locally inertial frames. Mass, previously postulated as closure residual $\Delta(\psi) = d(\psi, \Mcl)$, is reinterpreted as measuring a state's displacement from the manifold on which geodesic-like motion is defined.

This paper does not derive Einstein's field equations, identify $\Mcl$ with physical spacetime, compute gravitational coupling constants, or claim a completed theory of gravity. It establishes a controlled structural analogy between constraint-manifold geometry and gravitational phenomenology, providing a pathway --- not a destination --- for the closure framework's engagement with gravitational physics.
\end{abstract}

%==================================================================
\section{Introduction}\label{sec:intro}
%==================================================================

Paper~IX~\citep{SmartIX} established a formal continuous-limit description of the closure framework in which:
\begin{itemize}
    \item the closed state set $\Scl$ is represented as a constraint manifold $\Mcl$,
    \item the ambient inner product induces a Riemannian metric $g_{\Mcl}$ on the smooth strata of $\Mcl$,
    \item curvature of $\Mcl$ arises from the nonlinear interaction of constraint classes.
\end{itemize}

The present paper asks: \textit{can the curvature of the constraint manifold be interpreted as a structural analogue of gravitational influence?}

In general relativity, free particles follow geodesics of spacetime, and the curvature of spacetime is determined by the distribution of mass-energy. In the present framework, we propose an analogous picture: closure-compatible states evolve along geodesics of $\Mcl$, and the curvature of $\Mcl$ --- determined by the constraint structure, not by an external stress-energy tensor --- governs the geometry of admissible motion.

\subsection*{Scope and epistemic status}

This paper develops a gravity-facing \textit{interpretation} of the closure framework in which gravitational structure is associated with the curvature of the constraint manifold rather than assumed a priori as spacetime geometry. It does not derive Einstein's field equations, identify $\Mcl$ with physical spacetime, compute the Newton constant, or claim a completed theory of gravity. The analogy is structural, not dynamical: it shows that the framework admits a gravity-like geometric language, not that it reproduces general relativity.

%==================================================================
\section{Geometric Framework Recap}\label{sec:recap}
%==================================================================

From Paper~IX~\citep{SmartIX}:

\begin{itemize}
    \item \textbf{Ambient space:} $V_\infty = L^2(\Omega, \mu)$, a model space for the continuous-limit formalism.
    \item \textbf{Admissible domain:} $\Sadm \subset V_\infty$, normalised states satisfying boundary conditions.
    \item \textbf{Constraint manifold:} $\Mcl = \{\psi \in \Sadm \mid \Fcal(\psi) = 0\}$, the zero set of the closure functional. Used in the extended sense of a stratified constraint space with smooth manifold structure on regular strata.
    \item \textbf{Induced metric:} on smooth strata, $g_{\Mcl}(\xi, \eta) = \langle\xi, \eta\rangle$ for $\xi, \eta \in T_\psi \Mcl$.
    \item \textbf{Closure operator:} $\Ccal(\psi) \in \arg\inf_{\phi \in \Mcl} \|\psi - \phi\|$, metric projection onto $\Mcl$.
    \item \textbf{Closure residual:} $\Delta(\psi) = \|\psi - \Ccal(\psi)\| = d(\psi, \Mcl)$.
\end{itemize}

%==================================================================
\section{Constrained Motion on the Manifold}\label{sec:motion}
%==================================================================

\subsection{Curves on $\Mcl$}

\begin{definition}[Constrained curve]\label{def:curve}
A constrained curve is a smooth map $\gamma : I \to \Mcl$ (where $I \subset \mathbb{R}$ is an interval) such that $\gamma(t)$ lies on the constraint manifold for all $t \in I$.
\end{definition}

The tangent vector $\dot{\gamma}(t) \in T_{\gamma(t)} \Mcl$ at each point is required to lie in the tangent space of the manifold. This means the evolution is \textit{constrained}: only directions compatible with the closure conditions are admissible.

\subsection{Geodesic-like evolution}

\begin{definition}[Geodesic on $\Mcl$]\label{def:geodesic}
A geodesic on $\Mcl$ is a curve $\gamma : I \to \Mcl$ that extremises the arc-length functional with respect to the induced metric:
\begin{equation}\label{eq:geodesic}
    \delta \int_I \sqrt{g_{\Mcl}(\dot{\gamma}, \dot{\gamma})}\, dt = 0
\end{equation}
Equivalently, geodesics satisfy the geodesic equation with respect to the Levi-Civita connection of $g_{\Mcl}$:
\begin{equation}\label{eq:geoeq}
    \nabla_{\dot{\gamma}} \dot{\gamma} = 0
\end{equation}
where $\nabla$ is the connection induced by $g_{\Mcl}$.
\end{definition}

\begin{remark}
This is a geometric definition on the constraint manifold. The parameter $t$ is an arc-length or affine parameter along the curve, not identified with physical time. The geodesic equation~\eqref{eq:geoeq} is the standard differential-geometric geodesic equation, applied here to the constraint manifold rather than to a spacetime manifold.
\end{remark}

%==================================================================
\section{Curvature as Effective Gravitational Structure}\label{sec:gravity}
%==================================================================

\subsection{The gravitational analogy}

In general relativity:
\begin{itemize}
    \item free particles follow spacetime geodesics,
    \item spacetime curvature is sourced by mass-energy (via the Einstein equations),
    \item nearby geodesics converge or diverge according to the Riemann curvature tensor.
\end{itemize}

In the constraint-manifold framework:
\begin{itemize}
    \item closure-compatible states follow geodesics of $\Mcl$,
    \item manifold curvature arises from constraint interaction (not from an external source term),
    \item nearby geodesics on $\Mcl$ converge or diverge according to the curvature of the constraint surface.
\end{itemize}

\begin{remark}[Structural analogy, not derivation]
The analogy is structural: the mathematical form (geodesic motion on a curved manifold) is the same, but the physical content differs. In GR, the manifold is spacetime and curvature is sourced by stress-energy. In the present framework, the manifold is a constraint surface in state space and curvature arises from constraint interaction. Whether this analogy can be elevated to a physical correspondence is an open question.
\end{remark}

\subsection{Constraint-induced geodesic deviation}

On a curved manifold, nearby geodesics separate at a rate governed by the curvature. This is the geodesic deviation equation:
\begin{equation}\label{eq:deviation}
    \frac{D^2 \xi^\mu}{ds^2} = -R^\mu{}_{\nu\rho\sigma}\, \dot{\gamma}^\nu\, \xi^\rho\, \dot{\gamma}^\sigma
\end{equation}
where $\xi$ is the separation vector between nearby geodesics and $R$ is the Riemann curvature tensor of $g_{\Mcl}$.

In regions of positive curvature, nearby closure-compatible trajectories converge; in regions of negative curvature, they diverge. This is the constraint-manifold analogue of tidal gravitational effects.

%==================================================================
\section{Mass, Inertia, and the Manifold}\label{sec:mass}
%==================================================================

The mass--residual postulate (Papers~V, VII, IX) states:
\begin{equation}
    m = \kappa\, \Delta(\psi) = \kappa\, d(\psi, \Mcl)
\end{equation}

In the present geometric context, this acquires an additional interpretation:

\begin{itemize}
    \item \textbf{Mass as displacement:} mass measures how far a state is from the constraint manifold. Massless states lie on $\Mcl$; massive states are displaced from it.
    \item \textbf{Inertia as constraint resistance:} a state's resistance to changes in its closure-compatible trajectory is governed by its relationship to the manifold geometry. States farther from $\Mcl$ have larger closure residuals and, within the framework, greater effective mass.
    \item \textbf{Geodesic motion for massless states:} states on $\Mcl$ (with $\Delta = 0$) follow geodesics of the induced metric. States off $\Mcl$ are subject to a restoring tendency toward the manifold (the closure projection), in addition to any geodesic-like motion along it.
\end{itemize}

\begin{remark}
The statement ``mass appears as resistance to full closure, while curvature governs the geometry of admissible motion'' is an interpretive summary within the framework. It is not a derivation of the equivalence principle or of the inertial-gravitational mass identity.
\end{remark}

%==================================================================
\section{Effective Potential Picture}\label{sec:potential}
%==================================================================

The combined effect of closure residual and manifold curvature can be captured in a formal effective potential:
\begin{equation}\label{eq:potential}
    V_{\mathrm{eff}}(\psi) = \Delta(\psi)^2 + \lambda\, \mathcal{K}(\psi)
\end{equation}
where:
\begin{itemize}
    \item $\Delta(\psi)^2 = d(\psi, \Mcl)^2$: the squared closure residual, penalising displacement from the manifold.
    \item $\mathcal{K}(\psi)$: a local curvature proxy on $\Mcl$ (e.g.\ the scalar curvature or the norm of the second fundamental form at the nearest point on $\Mcl$).
    \item $\lambda > 0$: a structural coupling parameter, not derived.
\end{itemize}

This effective potential is a formal object: it suggests how curvature and closure may jointly shape admissible configurations, but it is not claimed as a dynamical Lagrangian or Hamiltonian. Its role is to provide a variational language connecting the closure residual (mass) to the manifold geometry (gravity-like curvature).

%==================================================================
\section{Local Flatness Analogue}\label{sec:flatness}
%==================================================================

In general relativity, the equivalence principle guarantees that any sufficiently small region of spacetime is approximately flat: local physics is approximately Minkowskian.

On the constraint manifold $\Mcl$, an analogous statement holds:

\begin{remark}[Local flatness on $\Mcl$]
In regions of $\Mcl$ where curvature is small (i.e.\ where constraint classes interact approximately linearly), the induced metric is approximately Euclidean and constrained motion is approximately linear. This provides the constraint-manifold analogue of a locally inertial frame.

In strongly curved regions --- where multiple constraint classes interact nonlinearly, such as near composite-class boundaries --- trajectories deviate from linearity and the effective gravitational analogue is strongest.
\end{remark}

%==================================================================
\section{Comparison with General Relativity}\label{sec:comparison}
%==================================================================

\begin{table}[ht]
\centering
\caption{Structural comparison between general relativity and the constraint-manifold framework.}
\label{tab:comparison}
\begin{tabular}{@{}ll@{}}
\toprule
General relativity & Constraint-manifold framework \\
\midrule
Spacetime manifold & Constraint manifold $\Mcl$ \\
Spacetime metric $g_{\mu\nu}$ & Induced metric $g_{\Mcl}$ \\
Geodesic motion of free particles & Geodesic-like motion of closure-compatible states \\
Curvature sourced by stress-energy & Curvature from constraint interaction \\
Mass-energy shapes geometry & Closure structure shapes manifold geometry \\
Equivalence principle (local flatness) & Approximate local flatness on smooth strata \\
Einstein field equations & Not derived; no analogue established \\
Gravitational waves & Not addressed \\
Black hole solutions & Not addressed \\
\bottomrule
\end{tabular}
\end{table}

\textsc{Status:} The entries in the left column are established physics. The entries in the right column are structural analogies within the present framework. The correspondence is at the level of geometric language, not at the level of dynamical equations. No Einstein field equations, stress-energy coupling, or gravitational constant are derived.

%==================================================================
\section{What Is and Is Not Claimed}\label{sec:claims}
%==================================================================

\subsection*{Claimed (within the framework)}

\begin{itemize}
    \item A gravity-facing geometric analogy exists: curvature of $\Mcl$ plays a structural role analogous to gravitational curvature.
    \item Constrained geodesic-like motion on $\Mcl$ is well-defined on smooth strata.
    \item Geodesic deviation on $\Mcl$ provides a structural analogue of tidal effects.
    \item Local flatness and trajectory bending have natural analogues on $\Mcl$.
    \item Mass (closure residual) and curvature (constraint interaction) jointly characterise the geometry within which admissible motion occurs.
\end{itemize}

\subsection*{Not claimed}

\begin{itemize}
    \item Exact equivalence to general relativity.
    \item Identification of $\Mcl$ with physical spacetime.
    \item Derivation of Einstein's field equations.
    \item Computation of the Newton constant or gravitational coupling.
    \item Gravitational waves, black holes, or cosmological solutions.
    \item The equivalence principle as a theorem of the framework.
    \item Experimental predictions distinguishing this framework from GR.
\end{itemize}

%==================================================================
\section{Limitations}\label{sec:limits}
%==================================================================

\begin{itemize}
    \item \textbf{No dynamical action or field equation.} The framework provides geometric structure but no evolution law. Constrained geodesic motion is defined, but the source of curvature is not coupled to a dynamical equation analogous to Einstein's.
    \item \textbf{No identification of physical time.} The curve parameter $t$ is geometric, not identified with physical coordinate time or proper time.
    \item \textbf{No stress-energy tensor analogue.} Curvature arises from constraint interaction, but no tensor equation relating curvature to a source term is established.
    \item \textbf{No equivalence principle proof.} Local flatness on smooth strata is a geometric observation, not a derivation of the equivalence principle.
    \item \textbf{No Newtonian limit.} The framework does not produce Newtonian gravity or its post-Newtonian corrections.
    \item \textbf{No coupling constants derived.} The structural parameter $\lambda$ in the effective potential and the mass--residual constant $\kappa$ are not computed.
    \item \textbf{No experimental predictions.} The framework does not yet produce predictions distinguishable from general relativity or competing theories.
    \item \textbf{$\Mcl$ is not spacetime.} The constraint manifold lives in state space, not in a spacetime of the kind used in general relativity. Whether and how the two can be related is an open question.
\end{itemize}

%==================================================================
\section{Conclusion}\label{sec:conclusion}
%==================================================================

Paper~IX introduced the constraint manifold $\Mcl$ and its induced metric. The present paper develops the first gravitational interpretation: curvature of $\Mcl$, arising from constraint interaction, provides a structural analogue of gravitational influence. Constrained geodesic-like motion on $\Mcl$ follows the induced metric; geodesic deviation provides an analogue of tidal effects; and approximate local flatness on smooth strata provides an analogue of locally inertial frames.

Mass (the closure residual $\Delta$) measures displacement from the manifold on which geodesic-like motion is defined. Curvature governs the geometry of that motion. Together, they provide a geometric language in which the closure framework can engage with gravitational phenomenology.

This is a pathway, not a destination. The analogy is structural: the geometric form matches, but the dynamical content of general relativity --- field equations, coupling constants, cosmological solutions --- is not reproduced. Whether the constraint-manifold framework can be developed into a dynamical theory of gravity, or whether it provides only a structural pre-geometric language, remains the central open question.

%==================================================================
\bibliographystyle{plainnat}
\bibliography{references}

\end{document}
